\section{Code--Data Identification and Data-Preserving Rewriting}
\label{sec:code_data}
\subsection{Why Data Identification Comes Before Transformation}
A raw binary contains no explicit distinction between instructions and
embedded strings or tables. String bytes may decode as plausible
instructions, so rewriting their apparent instruction sequence can corrupt
data. Table~\ref{tab:broken-payload} illustrates this failure; in
Table~\ref{tab:functional-payload}, the string instead occupies an
instruction operand that must remain intact.

Preparation separates two responsibilities. \emph{Identification} proposes
data regions; \emph{data handling} preserves supported objects and updates
their references after layout changes. A reference is an instruction
operand or address calculation pointing to code or data. Preserving a
string is insufficient if the program still reads its old address.
Figure~\ref{fig:extended_pipeline} shows the alternative symbolic/heuristic
and LLM classifiers feeding shared handling. Independently, instructions
can be transformed natively or by an LLM: classification does not imply
LLM transformation.
\begin{figure}[t]
\centering
\begin{tikzpicture}[
  box/.style={draw,rounded corners=2pt,align=center,font=\small,
    minimum height=0.8cm,text width=4.0cm},
  arrow/.style={-{Latex[length=2mm]},thick}]
\node[box,text width=5.7cm] (input) at (0,0) {Raw binary and target architecture};
\node[box] (sym) at (-2.35,-1.3) {Symbolic exploration\\and static heuristics};
\node[box] (llm) at (2.35,-1.3) {LLM data classifier\\(proposes regions only)};
\node[box,text width=7.4cm] (data) at (0,-2.65) {Shared data handling: check candidate regions,\\preserve supported data and track references};
\node[box,text width=7.4cm] (trans) at (0,-4.0) {Transform code: native methods, assembly wrapping,\\or LLM-proposed EQ, NOP, BCF, or IBP};
\node[box,text width=7.4cm] (build) at (0,-5.3) {Finalize layout and references; check construction};
\node[box,text width=7.4cm] (test) at (0,-6.6) {Test the produced binary with an external harness};
\node[box,text width=1.4cm,font=\scriptsize] (stop) at (5.4,-4.0) {Stop:\\no binary};
\draw[arrow] (input) -- (sym);
\draw[arrow] (input) -- (llm);
\draw[arrow] (sym) -- (data);
\draw[arrow] (llm) -- (data);
\draw[arrow] (data) -- (trans);
\draw[arrow] (trans) -- (build);
\draw[arrow] (build) -- node[right,font=\scriptsize] {accepted} (test);
\draw[arrow] (data.east) -- node[above,font=\scriptsize] {unsupported} (5.4,-2.65) -- (stop.north);
\draw[arrow] (build.east) -- node[below,font=\scriptsize] {invalid} (5.4,-5.3) -- (stop.south);
\node[font=\small,fill=white] at (0,-1.3) {or};
\end{tikzpicture}
\caption{Organization of the extensions. Identification proposes data
regions; shared handling and construction checks determine whether a
binary reaches the independent functional test.}
\label{fig:extended_pipeline}
\end{figure}

\subsection{What a Classification Result Means}
For an input of $n$ bytes, a proposed data region is an interval $[s,e)$:
it starts at offset $s$ from the beginning of the input and includes every
byte up to, but not including, $e$. For example, $[2,8)$ denotes six bytes.
An optional type describes a string encoding, such as ASCII or UTF-16LE,
or another data object. The output is a set of candidate intervals,
not a replacement binary.

Three stages must be distinguished: a classifier's proposed regions,
regions accepted by the engine's checks, and regions actually extracted
from the instruction stream. A correct proposal may still be unsupported
if its references cannot be updated. Conversely, different proposals can
produce the same binary when the shared checks retain the same objects.
The experiments compare annotations, accepted outputs, and final hashes
at their respective stages.

\subsection{Symbolic Exploration Combined with Static Heuristics}
The non-LLM method combines three sources of evidence. First, static
analysis recovers a control-flow graph (CFG): nodes represent basic blocks
of instructions and edges represent possible transfers of execution.
Zone uses angr's CFGFast analysis for this initial structure~\cite{angr_cfg_docs}.
Second, symbolic execution explores possible paths while representing
unknown register and memory values as symbolic expressions. This supplies
additional evidence about which instructions can be reached. Third,
pattern-based rules recognize printable ASCII and UTF-16LE string
candidates, including their terminators, and consider regions identified
as data by the control-flow analysis. Candidates overlapping confirmed
instruction bytes are filtered out. These rules supply recognizable
patterns and supporting evidence; an unexplored region alone is not
sufficient to classify it as data.

An observed instruction protects its entire decoded byte span, not only
its starting address. This matters for variable-length instructions: an
interior byte should not become a data candidate merely because the
explorer recorded only the instruction's first byte. Newly observed entry
points can also support another graph-recovery pass. Printable candidates
that overlap observed code are excluded by the string heuristic; candidates
overlapping graph-recovered code need additional structural support.

Exploration is bounded by time, solver, and state limits. The initial state
does not reproduce every feature of a Windows process, including all
external API behavior. Consequently, an unvisited region is \emph{unknown},
not proven to be data. Reports retain the evidence category, stopping
reason, unresolved states, and errors. These distinctions prevent a
partial execution search from being presented as a complete partition of
the binary into code and data.

\subsection{LLM Identification and Shared Data Handling}
The LLM receives hexadecimal bytes, a linear disassembly, and candidate
printable strings, and returns structured data intervals. The disassembly
is evidence to assess, not a trusted annotation. This mode neither
extracts data nor rewrites instructions.

Deterministic parsing rejects empty or out-of-bounds intervals and invalid
overlaps. The engine examines supported references and retains accepted
data separately from the instructions. After transformation, it finalizes
references using the new layout. Integrity checks compare retained bytes
with their saved originals at the final locations and verify supported
references against their expected destinations; mismatches reject the
construction. Address handling is architecture-specific.

For example, a benign routine may copy embedded strings \texttt{alpha}
and \texttt{beta} into a buffer. Identification locates the strings;
handling preserves their contents and links; the final execution test
checks the expected buffer. These test different stages of the pipeline.

Correction requests append the rejected answer and validator diagnostic
to the conversation and ask for a revised proposal, within the configured
attempt limit. An interval beyond the input, for example, receives an
out-of-bounds diagnostic. Classification feedback concerns data locations;
transformation feedback concerns assembly. Likewise, examples illustrating
classification are independent of the transformation zero-shot/two-shot
conditions in the next section.

Caching reuses identified regions associated with an input identity.
References are computed again for the working representation because
rewriting can move code and data. Reusing identification therefore does
not reuse obsolete addresses. Both classifiers feed the same handling,
which can reconcile different proposals into identical final binaries.
